\par\Needspace{6\baselineskip}\medskip
\hypertarget{native-110-article081}{}\pdfbookmark[2]{81.}{article-081}
\begin{center}81.\end{center}
Ist die Exzentrizität $e$ gegeben, so findet man den Winkel $\Pi$ aus der Gleichung
\[
\cos(A-\Pi)=-\frac{\cos(A-\tfrac12N-\tfrac12N')}{e\cos\tfrac12(N'-N)}
\]
nachdem man den Hilfswinkel $A$ aus der Gleichung
\[
\operatorname{tang}(A-\tfrac12N-\tfrac12N')=\frac{r'+r}{r'-r}\operatorname{tang}\tfrac12(N'-N)
\]
bestimmt hat. Die Mehrdeutigkeit, die bei der Bestimmung des Winkels $A-\Pi$ aus seinem Kosinus verbleibt, liegt in der Natur der Aufgabe, die daher zwei verschiedene Lösungen zulassen kann. Welche anzunehmen und welche zu verwerfen ist, muss anhand anderer Angaben entschieden werden. Dazu muss bereits wenigstens ein Näherungswert \SourceLeaf{111}für $\Pi$ bekannt sein. — Sobald $\Pi$ gefunden ist, berechnet man $p$ entweder nach den Formeln
\[
p=r(1+e\cos(N-\Pi))=r'(1+e\cos(N'-\Pi))
\]
oder nach der Formel
\[
p=\frac{2rr'e\sin\tfrac12(N'-N)\sin(\tfrac12N+\tfrac12N'-\Pi)}{r'-r}\mathpunct{\cdot}
\]
